\section{SMC${}^2$ over parameters}\label{sec:smc2}

\subsection{The tempered posterior over $\thetadyn$}

The bootstrap filter of Section~\ref{sec:pf} assumes the static parameter
$\thetadyn$ is known. In practice $\thetadyn$ is itself the principal object
of inference, with a prior $p(\thetadyn)$ and posterior
\begin{equation}
  p(\thetadyn \mid Y_{1:T})
    \;\propto\; p(\thetadyn)\;p(Y_{1:T}\mid\thetadyn).
\end{equation}
A direct sampler is unavailable because $p(Y_{1:T}\mid\thetadyn)$ has no
closed form. The SMC${}^2$ algorithm of Chopin, Jacob and Papaspiliopoulos
(2013) replaces this likelihood by the unbiased estimator
$\Lhat_N(\thetadyn)$ from \eqref{eq:unbiased-likelihood} and, for a fixed
schedule $0 = \lambda_0 < \lambda_1 < \cdots < \lambda_K = 1$, defines a
sequence of tempered targets
\begin{equation}
  \pi_\lambda(\thetadyn) \;\propto\;
    p(\thetadyn)^{1-\lambda}\,\Lhat_N(\thetadyn)^{\lambda},
    \qquad \lambda \in [0,1].
  \label{eq:tempered-target}
\end{equation}
At $\lambda=0$ the target is the prior; at $\lambda=1$ it is, in
expectation, the posterior. The interpolation in between bridges the two
in a way that keeps the importance weights well-behaved.

\subsection{Adaptive tempering and the move step}

A population of $M$ \emph{parameter particles}
$\{\thetadyn^{(m)}\}_{m=1,\dots,M}$ is propagated through the schedule
$\{\lambda_k\}$. At each step the unnormalised tempering weight
\begin{equation}
  \omega_k^{(m)} \;=\;
    \biggl(\frac{\Lhat_N(\thetadyn^{(m)})}{p(\thetadyn^{(m)})^{-1}}\biggr)
    ^{\lambda_k - \lambda_{k-1}}
\end{equation}
is computed, and an effective sample size
$M_{\mathrm{eff}} = (\sum_m \omega_k^{(m)})^2 / \sum_m (\omega_k^{(m)})^2$
is monitored. The interval $\lambda_k - \lambda_{k-1}$ is chosen
\emph{adaptively}, by bisection, to drive $M_{\mathrm{eff}}/M$ to a
prescribed target (typically $0.5$); this avoids both unnecessarily
short bridges and weight collapse. After resampling, each parameter
particle is moved by a Hamiltonian Monte Carlo (HMC) kernel that targets
$\pi_{\lambda_k}$, with the mass matrix re-estimated from the current
particle cloud. Algorithm~\ref{alg:smc2} summarises the procedure.

\begin{algorithm}[H]
\caption{Adaptive-tempered SMC${}^2$ over $\thetadyn$}
\label{alg:smc2}
\begin{algorithmic}[1]
\State Sample $\thetadyn^{(m)} \sim p(\thetadyn)$ for $m=1,\dots,M$;
       set $\lambda_0 = 0$.
\State For each $m$, run Algorithm~\ref{alg:sir} with $\thetadyn^{(m)}$
       to obtain $\Lhat_N(\thetadyn^{(m)})$.
\While{$\lambda_k < 1$}
  \State Solve for the smallest increment $\delta \in (0,1-\lambda_k]$
         such that the tempering ESS ratio with weights
         $\omega^{(m)} \propto \Lhat_N(\thetadyn^{(m)})^{\delta}$
         equals the target (e.g.\ $0.5$).
         Set $\lambda_{k+1} = \lambda_k + \delta$.
  \State Resample $\{\thetadyn^{(m)}\}$ systematically by
         $\{\omega^{(m)}\}$.
  \State \textbf{Move.} Run a small number of HMC steps targeting
         $\pi_{\lambda_{k+1}}$ on each $\thetadyn^{(m)}$, re-running
         Algorithm~\ref{alg:sir} as needed to evaluate the unbiased
         likelihood at proposed values.
  \State $k \leftarrow k+1$.
\EndWhile
\end{algorithmic}
\end{algorithm}

\subsection{Rolling-window SMC${}^2$ and the Schr\"odinger--F\"ollmer bridge}

When data arrive sequentially, running Algorithm~\ref{alg:smc2} from
scratch over the full record at every new arrival is wasteful. The
implementation in this code base instead operates over a sliding sequence
of \emph{windows}: with a fifteen-minute time grid, each window is one
day long ($96$ bins), the stride between consecutive windows is twelve
hours ($48$ bins), and the macrocycle of fourteen days yields
$(14 \cdot 96 - 96)/48 + 1 = 27$ windows.

For window $w \ge 2$ the parameter cloud is \emph{warm-started} from the
posterior of window $w-1$. The transition is mediated by a
Schr\"odinger--F\"ollmer bridge between two Gaussian fits, parameterised
along the Bures--Wasserstein geodesic in the manifold of positive-definite
covariance matrices. Concretely, let $(m_0,\Sigma_0)$ denote a Gaussian
moment match to the previous window's posterior, and let
$(m_1,\Sigma_1)$ denote an importance-weighted moment match obtained by
a single particle-filter evaluation per particle on the new window's data.
For a blend parameter $\rho \in [0,1]$ the bridge measure has mean and
covariance
\begin{align}
  m_\rho &= (1-\rho)\,m_0 + \rho\,m_1, \\
  \Sigma_\rho &= \bigl((1-\rho)I + \rho\,T\bigr)\,\Sigma_0\,
                 \bigl((1-\rho)I + \rho\,T\bigr)^{\!\top},
\end{align}
where $T = \Sigma_0^{-1/2}\bigl(\Sigma_0^{1/2}\Sigma_1\Sigma_0^{1/2}\bigr)^{1/2}\Sigma_0^{-1/2}$
is the Bures--Wasserstein optimal-transport map. The default
implementation uses $\rho = 0.7$, biasing the warm-start toward the new
window's evidence. We do not develop the geometry here; the reader
unfamiliar with it can treat the bridge as a black box that produces a
sensible Gaussian initialisation for the next window's tempering schedule.
